
\subsubsection{4.1 Research Questions}

\textbf{RQ1 (Orbit Existence, H-E1):} Do scaling and sign-flip symmetry orbits have non-negligible geometric diameter in the Schürholt MNIST MLP zoo, and how large are they?

\textbf{RQ2 (Encoder Invariance, H-M1):} Is NFT, trained on raw Schürholt zoo weights, naturally invariant to scaling and sign-flip orbits?

\textbf{RQ3 (Geometric Concentration, H-M2):} Does scaling canonicalization concentrate geometric structure --- specifically, does it improve PCA explained variance and linear-probe property prediction?

\textbf{RQ4 (Uniqueness Audit, H-C1):} Is sign-flip canonicalization via the majority-sign algorithm well-defined for the Schürholt MNIST MLP architecture?

An implicit RQ5 (H-M3) evaluates end-to-end property prediction improvement under canonicalization, providing directional evidence and establishing scale requirements.

\subsubsection{4.2 Dataset}

\textbf{Schürholt MNIST Model Zoo} [Schürholt et al., 2022]. 2-layer MLPs ($784\rightarrow 64\rightarrow 10$, ReLU activations) trained on MNIST with varying learning rates (log-uniform in [$10^{-5}, 10^{-1}])$, weight decay values, and initialization seeds. Ground-truth property labels:

\begin{table}[htbp]
\centering
\resizebox{\columnwidth}{!}{%
\begin{tabular}{lll}
\toprule
Property & Description & Approximate Range \\
\midrule
Test accuracy & MNIST test set accuracy & 0.60 -- 0.98 \\
Generalization gap & Train minus test accuracy & 0.00 -- 0.35 \\
Learning rate & Training LR used (log-scale) & $10^{-5}$ -- $10^{-1}$ \\
\bottomrule
\end{tabular}
}
\end{table}

We use a locally archived subset of N=500 models. For orbit characterization (RQ1, RQ2), this subset is sufficient --- relative comparisons produce tight CIs at N=500. For property prediction (RQ3, RQ5), N=500 is severely underpowered, as we quantify in Section 5.

\subsubsection{4.3 Baselines}

\textbf{Raw weights $\rightarrow$ NFT (Condition A).} Standard NFT without preprocessing. At N=500, NFT achieves Spearman $\rho \approx$ 0 across all tasks (all CI-covered values are statistically indistinguishable from 0); see Section 5.5.

\textbf{Random normalization $\rightarrow$ NFT (Condition E).} Null control that applies semantically meaningless normalization.

\textbf{Layer statistics} \citep{Unterthiner2020PredictingNN}. Layer-wise weight distribution moments. Achieves Spearman $\rho \approx$ 0.9 on simple zoos. Included to contextualize the NFT gap at N=500.

\subsubsection{4.4 Evaluation Metrics}

\textbf{Spearman $\rho$.} Rank correlation between predicted and ground-truth property labels on the 50-model held-out test split.

\textbf{Invariance gap.} mean(within-orbit similarity) $-$ mean(cross-orbit same-property similarity), with bootstrap 95\% CI.

\textbf{Explained Variance Ratio (EVR).} Fraction of total weight-matrix variance captured by top-k PCs.

\textbf{Linear regression $R^2$.} Out-of-sample $R^2$ from top-k PCs onto property labels.

\textbf{Fraction unique.} Fraction of zoo models for which majority-sign produces a unique canonical form.

\subsubsection{4.5 Implementation Details}

\textbf{NFT.} d\_model=256, 4 self-attention layers, 8 heads, CLS-token pooling. Weights tokenized row-by-row from each weight matrix (one token per neuron incoming-weight vector). NFT has approximately 3.4M parameters.

\textbf{Training.} Adam, lr=1e-3 (primary) or lr=$3\times 10^{-4}$ (fallback), weight\_decay=1e-4. Maximum 50 epochs with early stopping on validation Spearman $\rho$ (patience=10).

\textbf{Data split.} 80/10/10 fixed: 400 training, 50 validation, 50 test. Same split across all conditions.

\textbf{Oracle orbit construction.} $\alpha _i$ ∼ log-Uniform(0.1, 10.0) for scaling; s\_i ∼ Bernoulli(0.5) $\rightarrow$ {-1, +1} for sign-flip. Seeds fixed for reproducibility.

\textbf{Cross-orbit pairs (H-M1).} For each base model, cross-orbit partner is sampled as the nearest model in the same test-accuracy decile.

\textbf{Bootstrap CI.} scipy.stats.bootstrap with n\_boot=1,000, seed=42, 95\% confidence level.

\textbf{Hardware.} All experiments run on CPU (NFT training) or GPU (H-M1 checkpoint training). NFT training takes approximately 5--15 minutes per condition for N=400 training samples.

